%%%%%%%%%%%%%%%%%
% This is an sample CV template created using altacv.cls
% (v1.7.4, 30 July 2025) written by LianTze Lim (liantze@gmail.com). Compiles with pdfLaTeX, XeLaTeX and LuaLaTeX.
%
%% It may be distributed and/or modified under the
%% conditions of the LaTeX Project Public License, either version 1.3
%% of this license or (at your option) any later version.
%% The latest version of this license is in
%%    http://www.latex-project.org/lppl.txt
%% and version 1.3 or later is part of all distributions of LaTeX
%% version 2003/12/01 or later.
%%%%%%%%%%%%%%%%

\documentclass[10pt,a4paper,withhyper]{altacv}

% Layout and typography
\geometry{left=1.35cm,right=1.35cm,top=1.35cm,bottom=1.35cm}

\iftutex
  \setmainfont{Roboto Slab}
  \setsansfont{Lato}
  \renewcommand{\familydefault}{\sfdefault}
\else
  \usepackage[rm]{roboto}
  \usepackage[defaultsans]{lato}
  \renewcommand{\familydefault}{\sfdefault}
\fi
\usepackage{microtype}
\DisableLigatures{encoding = *, family = *}

% Restrained aerospace-inspired palette
\definecolor{ResumeText}{HTML}{111827}
\definecolor{ResumeMuted}{HTML}{4B5563}
\definecolor{ResumeAccent}{HTML}{163A5F}
\definecolor{ResumeRule}{HTML}{CBD5E1}
\colorlet{name}{ResumeText}
\colorlet{tagline}{ResumeAccent}
\colorlet{heading}{ResumeAccent}
\colorlet{headingrule}{ResumeRule}
\colorlet{subheading}{ResumeText}
\colorlet{accent}{ResumeAccent}
\colorlet{emphasis}{ResumeText}
\colorlet{body}{ResumeMuted}

\renewcommand{\namefont}{\Huge\bfseries}
\renewcommand{\personalinfofont}{\footnotesize}
\renewcommand{\cvsectionfont}{\large\bfseries}
\renewcommand{\cvsubsectionfont}{\normalsize\bfseries}
\renewcommand{\cvItemMarker}{{\small\textbullet}}

\setlength{\parindent}{0pt}
\setlength{\parskip}{0pt}
\AtBeginEnvironment{itemize}{\small}

% Local helpers for a one-column ATS-friendly layout
\newcommand{\contactsep}{\enspace\textbar\enspace}
\newcommand{\resumesection}[1]{%
  \cvsection{#1}\vspace{-0.45\baselineskip}%
}
\newcommand{\resumeentry}[4]{%
  \textbf{#1}%
  \ifstrempty{#3}{}{ {\small\color{emphasis}(#3)}}%
  \par
  \ifstrempty{#2}{}{%
    {\emph{#2}}%
    \ifstrempty{#4}{}{%
      \hfill{\small\color{emphasis}#4}%
    }%
    \par
  }%
}
\newcommand{\skillgroup}[2]{%
  \textbf{#1}: #2\par
}

\begin{document}

% TODO: Confirm whether the candidate wants to keep the full Lattes URL in the header or move it to a shorter link label.
% TODO: Confirm whether the MBA should remain if the resume is later tightened further for one-page submissions.
% TODO: Confirm whether graph attention should remain attributed to the Loyal Wingmen project in all target applications, or only when the 2026 journal article is cited.

{\namefont\color{name} Davi Guanabara Arag\~ao\par}
{\large\bfseries\color{tagline} AI/ML Engineer \textbar{} PhD Researcher in Autonomous Systems and Data Science\par}
\vspace{0.35em}
{\small
S\~ao Paulo, Brazil\contactsep
\href{mailto:davi\_guanabara@live.com}{davi\_guanabara@live.com}\contactsep
+55 85 9 8673-6026\contactsep
\href{https://github.com/DaviGuanabara}{github.com/DaviGuanabara}\contactsep
\href{https://linkedin.com/in/davi-guanabara}{linkedin.com/in/davi-guanabara}\contactsep
\href{http://lattes.cnpq.br/2180140366316472}{lattes.cnpq.br/2180140366316472}\par
}
\vspace{0.55em}

\resumesection{Professional Summary}

AI/ML Engineer and PhD researcher at the Aeronautics Institute of Technology (ITA), applying artificial intelligence, machine learning, and software engineering to autonomous systems, geospatial computing, and aerospace applications. Currently contributes to research initiatives in collaboration with the Brazilian Airspace Control System (SAC/ICEA), developing computational methods for urban airspace planning, intelligent decision-making, and simulation-driven experimentation. Research interests focus on reinforcement learning, multi-agent systems, geospatial AI, and reproducible research software for complex engineering problems.

\medskip
\resumesection{Areas of Expertise}

\skillgroup{Autonomous Systems}{
Autonomous Decision-Making • Multi-Agent Systems • Autonomous UAV Systems
}

\skillgroup{Machine Learning}{
Deep Learning • Reinforcement Learning • Attention-Based Models
}

\skillgroup{Software Engineering}{
Modular Software Architecture • Event-Driven Systems • Experimentation Frameworks
}

\skillgroup{Geospatial Computing}{
Geospatial Modeling • Spatio-Temporal Modeling
}

\medskip
\resumesection{Selected Projects}

\resumeentry{Deep Reinforcement Learning with Graph Attention for Cooperative Loyal Wingmen}
{CTEDS Research Program | Brazil--Sweden Air Domain Study}
{2024--2026}{}

\begin{itemize}

\item Developed a deep reinforcement learning framework for cooperative threat-engagement decision-making in Loyal Wingmen UAVs operating under shared spatio-temporal observations.

\item Proposed a graph-attention neural architecture enabling adaptive information sharing and coordinated policy learning under partial observability.

\item Designed and evaluated simulation-based experiments integrating graph neural networks with multi-agent reinforcement learning for autonomous aerial systems.

\item Conducted within the Cooperative Threat Engagement with Heterogeneous Drone Swarms (CTEDS) research program, an international Brazil--Sweden collaboration involving government, academia, and industry.

\item Resulted in two peer-reviewed publications:

      \emph{Journal of Intelligent \& Robotic Systems} (Springer Nature, 2026, DOI:\href{https://doi.org/10.1007/s10846-026-02415-8}{10.1007/s10846-026-02415-8}) and \emph{2024 Brazilian Symposium on Robotics (SBR) and Workshop on Robotics in Education}(DOI:\href{https://doi.org/10.1109/SBR/WRE63066.2024.10837749}{10.1109/SBR/WRE63066.2024.10837749}).
\end{itemize}

\resumeentry{SkyWeaver --- Dynamic Airspace Configuration Platform for BRUTM}
{Research Project in collaboration with SAC/ICEA}
{2025--Present}{}

\begin{itemize}
\item Designed a reactive architecture supporting Dynamic Airspace Configuration to contribute with the Brazilian UTM (BRUTM) concept.
\item Developed modular software components for geospatial scene generation, operational state management, routing, and airspace infrastructure modeling.
\item Investigating graph-based optimization methods for droneport placement, route planning, and infrastructure provisioning under operational constraints.
\item Ongoing research explores software engineering, optimization, and AI techniques for future UTM, U-Space, and Advanced Air Mobility (AAM) operations.
\end{itemize}

\resumeentry{OndeVale --- A Property Valuation Platform}{Cloud-Deployed AI System}{2025}{}
\begin{itemize}
\item Developed a cloud-deployed platform for residential property valuation combining geospatial analytics, machine learning, and interactive web visualization.
\item Designed the complete machine learning pipeline, including data preprocessing, feature engineering, model training, inference services, and deployment.
\item Built a FastAPI backend and Vue.js frontend integrating GIS visualization with real-time model inference.
\item Deployed the application as a publicly accessible cloud service, connecting data preprocessing, model inference, a REST API, and an interactive geospatial interface.
\end{itemize}

\resumeentry{A Modular Framework for Reproducible Supervised Learning: A Vineyard Leaf Wetness Prediction Case Study}{}{2024--2025}{}
\begin{itemize}
\item Developed a deep learning model for vineyard leaf wetness prediction using meteorological observations from the Iowa Environmental Mesonet.
\item Designed a permutation-invariant Deep Sets architecture with masked attention to model heterogeneous spatio-temporal neighborhoods.
\item Built a reproducible supervised-learning framework integrating stochastic contextual sampling, Bayesian hyperparameter optimization, bootstrap evaluation, and deterministic seed management.
\end{itemize}

\medskip
\resumesection{Education}
\resumeentry{Ph.D.\ in Electronic Engineering and Computing}{Aeronautics Institute of Technology (ITA)}{2025--Present}{}
\resumeentry{Graduate Certificate in Data Science}{Aeronautics Institute of Technology (ITA)}{2024--2026}{}
\resumeentry{M.S.\ in Aeronautical Computing}{Aeronautics Institute of Technology(ITA)}{2021--2024}{}
\resumeentry{B.S.\ in Control and Automation Engineering}{University of Fortaleza (UNIFOR)}{2013--2019}{}



\end{document}
